% Pursuing Stacks (cotes 134-2 to 134-8): archive cote 134-6, site page 91 (PDF page 92), typescript page 409.
% typescript. Draft compared: ps5.tex:6127-6194.
% AI-transcribed by Claude Opus 5.5 (claude-opus-5-5, adaptive thinking, effort high), Message Batches API, batch ps3 (msgbatch_01XdFmuVVBAcgc7EFxdtCqsj), 2026-11-12; not yet checked by eye.
\scanpage{134-6/91}
\mspage{409}

structure in this abelianized duality context,  embodied by the \uline{tensor}
\uline{product structures} on both abelian topoi $A^{\wedge}_{ab}$ and $B^{\wedge}_{ab}$ ,  as contemplated
in section 104 C) ,  in a somewhat more general context.  Corresponding to
this extra structure on $A^{\wedge}_{ab}$ say,  we \xout{havex} saw that \xout{\unsure{$A^{\wedge}_{ab}$}xxx} this category
of abelian sheaves ``operates'' \inserted[typed]{covariantly} (by an operation denoted by \hand{$\circledast$}$_{\underline{Z}}$ or simply
\hand{$\circledast$} )  on any category of \hand{$\mathcal{M}$}-valued cosheaves on $A^{\wedge}_{ab}$ ,  and contravariantly
(by an operation  denoted by $\underline{\mathrm{Hom}}_{\underline{Z}}$ or simply \uline{Hom} , if no confusion may
arise)  on any category of \hand{$\mathcal{N}$}-valued  sheaves on $A^{\wedge}_{ab}$ ,  where \hand{$\mathcal{M},\mathcal{N}$} are
additive categories, stable under small direct resp. inverse limits.  The
\unsure{Latter|latter} operation (cf.  page 392)
\[ (65)\qquad (\over{\unsure{F}}{L},K)\hand{$\mapsto$}\underline{\mathrm{Hom}}_{\underline{Z}}(\over{\unsure{F}}{L},K) \ :\ A^{\wedge\,o}_{ab}\ \mathrm{x}\ A^{\wedge}_{\hand{$\mathcal{N}$}} \hand{$\longrightarrow$} A^{\wedge}_{\hand{$\mathcal{N}$}} \quad , \]
involving sheaves,  will be \xout{kept} used in the sequel ``tel quel'' ,  it's defi% NOTE: the right edge of the scan cuts this line; a hyphen after "defi" is not visible
\margin{$\times$} nition is of a tautological character,  independant of duality.  As for
the former operation involving cosheaves, we may view it via \xout{(59)} the
duality relation (59) as an operation of $A^{\wedge}_{ab}$ on \hand{$\mathcal{M}$}-valued \uline{sheaves} on the
dual abelian topos $B^{\wedge}_{ab}$ ,  and this operation will be denoted by the same
symbol \hand{$\circledast_{\underline{Z}}$} :
\margin{(L, M') [circled]}
\[ (66)\qquad (\over{\illegible{1}}{L},\over{\illegible{1}}{M'})\hand{\unsure{$\longrightarrow$|$\longmapsto$}} \over{\illegible{1}}{L}\hand{$\circledast_{\underline{Z}}$}\over{\illegible{1}}{M'} \ :\ A^{\wedge}_{ab}\ \mathrm{x}\ B^{\wedge}_{\hand{$\mathcal{M}$}} \hand{$\longrightarrow$} B^{\wedge}_{\hand{$\mathcal{M}$}} . \]
Replacing A by B in (66),  we get an operation of $B^{\wedge}_{ab}$  upon $A^{\wedge}_{\hand{$\mathcal{M}$}}$ ,  \xout{which}
\[ (67)\qquad (M,L')\hand{\unsure{$\to$}} M\hand{$\circledast_{\underline{Z}}$}L' \ :\ A^{\wedge}_{\hand{$\mathcal{M}$}}\ \mathrm{x}\ B^{\wedge}_{\xout{\illegible{1}}ab} \hand{\unsure{$\to$}} A^{\wedge} \quad . \]
Whenever convenient, we'll allow ourselves to write L' \hand{$\circledast$} M  instead of
M \hand{$\circledast$} L' (which does'nt seem to lead to any trouble),  and will henceforth
(unless special need should arise) drop the subscripts $\underline{Z}$ .

Thus,  for a given category  $A^{\wedge}$  of \hand{$\mathcal{M}$}-valued sheaves on $A^{\wedge}_{ab}$ ,  there
is a twofold operation on this category,  namely  $A^{\wedge}_{ab}$ itself operates
(the operation defined by L in $A^{\wedge}_{ab}$ depending contravariantly on L) as
well as the dual abelian topos $B^{\wedge}_{ab}$ (the operation defined by L' in $B^{\wedge}_{ab}$
depending covariantly on L'), \hand{$\mathcal{M}$} being any \xout{abelian} additive category
stable under small \xout{direct}\inserted[typed]{inverse} and \xout{inverse}\inserted[typed]{direct} limits (in order to insure existence
of both types of operations).  I would like to dwell a little more on this
twofold structure,  as I do'nt feel to have understood it thoroughly yet.
It is this second \xout{structure}\inserted[typed]{operation} mainly which \xout{isnt} has'nt become really fami-%
liar yet,  still less it's relationship to the first\hand{,} more familiar

\margin{91}
